\addcontentsline{toc}{section}{چکیده انگلیسی}
\thispagestyle{empty}

\begin{latin}
\begin{center}

{\huge Design and Implementation of a Multi-Agent Digital Twin Framework for Intelligent Monitoring, Prediction, and Decision-Making in Industrial Systems}

\vspace{1cm}

{\LARGE{Navid Nasiri}}

\vspace{0.2cm}

{\small navidnasyry@gmail.com}

\vspace{0.5cm}

2026

\vspace{0.5cm}

Department of Electrical and Computer Engineering

\vspace{0.2cm}

Isfahan University of Technology, Isfahan 84156-83111, Iran

\vspace{0.2cm}

Degree: M.Sc. \hspace*{3cm} Language: Farsi

\vspace{1cm}

{\small\textbf{Supervisor: Prof. Masoudreza Hashemi (hashemim@iut.ac.ir)}}

{\small\textbf{Advisor: Prof. Aref Karimi Afshar (a.karimiafshar@iut.ac.ir)}}
\end{center}
~\vfill


\noindent\textbf{Abstract}

\begin{small}
\baselineskip=0.6cm

With the increasing complexity of industrial equipment and the costs associated with unplanned downtime, the development of systems for continuous condition monitoring, failure prediction, and maintenance decision support has become essential. Despite the capability of digital twins to connect physical assets with their virtual representations, integrating Remaining Useful Life prediction, near-real-time monitoring, human--system interaction, and safe decision-making into extensible architectures remains a challenge. This study aims to design, implement, and evaluate a multi-agent digital twin framework for monitoring industrial equipment and predicting failures by leveraging deep learning and large language models. The proposed architecture consists of five layers: physical entity, data acquisition, data integration, digital twin agents, and application agents. To evaluate the framework, a digital twin of turbofan engines was developed using the FD001 subset of the NASA C-MAPSS dataset. This subset contains data from 100 engines for training and 100 engines for testing. To account for the operational risk of unexpected failures, Pinball loss with a quantile coefficient of 0.125 was employed to deliberately favor conservative underestimation of Remaining Useful Life. Compared with a baseline model based on a Long Short-Term Memory network trained using Mean Squared Error, the final model reduced the Mean Absolute Error from 17.56 to 17.21, increased the coefficient of determination from 0.24 to 0.71, and raised the underestimation ratio from 0.47 to 0.87. Independent monitoring, prediction, dashboard, and language-interaction agents were also implemented to enable near-real-time updates of the digital twin state, provision of relevant information, and execution of predefined operational commands. The results indicate that the proposed framework provides a modular, replaceable, and extensible foundation for safe, interactive, and adaptable predictive maintenance across other industrial assets and datasets.


\end{small}

\vspace{0.5 cm}

\noindent \textbf{Key Words}: Digital Twin, Multi-Agent System, Multi-Agent Digital Twin, Large Language Model, Predictive Maintenance, Remaining Useful Life, LSTM, NASA CMAPSS, Real-Time Monitoring, Explainable Artificial Intelligence.

\end{latin}